\documentclass[11pt,a4paper]{article}
\usepackage[margin=25mm]{geometry}
\usepackage{lmodern}
\usepackage{amsmath,amssymb,amsthm}
\usepackage{booktabs,microtype}
\usepackage[hidelinks]{hyperref}
\hypersetup{pdftitle={Analytic reductions for the Generalized Riemann Hypothesis},
  pdfauthor={GRH distillation project}}
\newtheorem{theorem}{Theorem}[section]
\newtheorem{lemma}[theorem]{Lemma}
\newtheorem{proposition}[theorem]{Proposition}
\newtheorem{corollary}[theorem]{Corollary}
\theoremstyle{definition}
\newtheorem{definition}[theorem]{Definition}
\theoremstyle{remark}
\newtheorem{remark}[theorem]{Remark}
\numberwithin{equation}{section}
\setcounter{tocdepth}{1}
\renewcommand{\Re}{\operatorname{Re}}
\renewcommand{\Im}{\operatorname{Im}}
\DeclareMathOperator{\Res}{Res}
\emergencystretch=1em
\title{Analytic reductions for the\ Generalized Riemann Hypothesis\\[0.4em]
  \large A mathematical distillation in progress}
\author{GRH distillation project}
\date{21 November 2026}

\begin{document}
\maketitle
\begin{abstract}
We develop explicit analytic reductions relevant to the announced proof of
the Generalized Riemann Hypothesis for Dirichlet $L$-functions. We give
complete proofs of the reduction to primitive characters, the normalized
functional equation, and several implications from cancellation estimates
to zero-free regions. We also examine which deductions from the quasi-Riemann
hypothesis methods can be reused, and prove precise limitations of a
particular second-moment argument and of nonuniform limiting arguments.
The decisive arithmetic estimates needed for a complete proof of GRH are
not established in this manuscript. Each conditional result identifies its
input explicitly. The purpose is to make the mathematical obligations
clear and to preserve the parts that can already be checked in ordinary
mathematical language.
\end{abstract}
\tableofcontents
\clearpage

\section{Purpose, scope, and conventions}
\label{sec:introduction}

The original work motivating this project is Levent Alp\"oge's announced
proof of the Riemann Hypothesis and the Generalized Riemann Hypothesis,
developed with Claude and accompanied by a released formal
source~\cite{Alpoge,AlpogeSource}. The project brief identifies the model as
Fable~5.2; the announcement itself names Claude without specifying a
model version. We attribute the announced proof to Alp\"oge and distinguish
it from the expository reconstructions and checks carried out here.

The assertion under study is GRH for Dirichlet $L$-functions. For a
Dirichlet character $\chi$, it states that
\begin{equation}\label{eq:grh-target}
 L(s,\chi)=0,\qquad 0<\Re s<1
 \quad\Longrightarrow\quad \Re s=\frac12.
\end{equation}
The interval $0<\Re s<1$ is the open critical strip. Definitions of the
characters and the functions are given in
Section~\ref{sec:primitive-reduction}. The principal character of conductor
one gives the Riemann zeta function. No assertion about all automorphic
$L$-functions is made here.

This manuscript does not yet prove~\eqref{eq:grh-target}. It proves
reductions whose missing arithmetic inputs are stated separately.
In particular, the functional equation reduces the critical-strip
assertion to nonvanishing for primitive characters to the right of
one half. We give two complementary ways to obtain nonvanishing from
estimates: a Mellin transform of smoothed M\"obius sums, and a positivity
argument for suitable weighted prime sums. These are implications;
the required cancellation and upper bounds remain substantial mathematical
tasks.

The quasi-RH sources~\cite{OpenAISeptember,OpenAIOctober} provide useful
comparisons. The September source states a zero-free half-plane
$\Re s>7/8$, and the October source isolates a simpler route to
$\Re s>11/12$. We prove the elementary consequences of the October
second-moment input and identify exactly what its exponent permits.
Neither stated half-plane is the GRH half-plane $\Re s>1/2$.
The automorphic estimates in those sources are not treated as having
been independently proved merely because their consequences are proved
here.

For positive real $x$, complex powers always mean
$x^s=\exp(s\log x)$ with the real logarithm. The notation
$F\ll_{\mathcal P}G$, where $G\geq0$, means that $|F|\leq CG$ for a
constant $C$ depending only on the listed fixed parameters $\mathcal P$.
Every parameter declared fixed is chosen before a limit such as
$X\to\infty$. Dependence on a fixed character is allowed unless a
uniform assertion over a character family is expressly made.

We use elementary complex analysis, finite Fourier orthogonality,
unique factorization of integers and ideals, and the ordinary integration
theorems for absolutely convergent sums. The analytic normalizations
needed below are included explicitly. The mathematical exposition is
organized by logical dependence. Source-file paths, build records, and
the continuing audit are maintained outside the paper.

\section{Two elementary analytic tools}
\label{sec:analytic-tools}

We record the Fourier and gamma-function facts used in the primitive
functional equation. Their usual normalizations can be compared with
\cite[\S\S1.8(iv), 5.2]{DLMF}.

\subsection{Poisson summation and the Gaussian}

A function $f:\mathbb R\to\mathbb C$ is a Schwartz function if it is
smooth and $\sup_x |x^j f^{(k)}(x)|<\infty$ for every pair of
nonnegative integers $j,k$. Define its Fourier transform by
\[
 \widehat f(y)=\int_{\mathbb R}f(x)e^{-2\pi ixy}\,dx.
\]
\begin{lemma}\label{tools:poisson}
For a Schwartz function,
\[
 \sum_{n\in\mathbb Z}f(n)=\sum_{m\in\mathbb Z}\widehat f(m).
\]
For $A>0$ the transforms of $e^{-\pi Ax^2}$ and
$xe^{-\pi Ax^2}$ are, respectively,
\[
 A^{-1/2}e^{-\pi y^2/A},\qquad
 -iA^{-3/2}y e^{-\pi y^2/A}.
\]
\end{lemma}
\begin{proof}
The periodization $F(x)=\sum_{n\in\mathbb Z}f(x+n)$ and all its
derivatives converge uniformly for $x\in[0,1]$. Its $m$th Fourier
coefficient is, by an absolutely convergent interchange,
\[
 \int_0^1F(x)e^{-2\pi imx}\,dx
  =\int_{\mathbb R}f(x)e^{-2\pi imx}\,dx=\widehat f(m).
\]
Integration by parts twice makes these coefficients absolutely summable.
The Fourier series of the smooth periodic function $F$ therefore
converges uniformly to $F$. Evaluating at zero proves the sum identity.

Let $G(y)$ be the transform of $e^{-\pi Ax^2}$. Differentiation under
the integral and integration by parts give
$G'(y)=-(2\pi y/A)G(y)$. The Gaussian integral, obtained by squaring
the integral and using polar coordinates, gives $G(0)=A^{-1/2}$.
Solving this differential equation proves the first formula. The transform
of $xe^{-\pi Ax^2}$ is $iG'(y)/(2\pi)$, proving the second.
All differentiations and boundary terms are justified by Gaussian decay.
\end{proof}

Applying the sum identity to $x\mapsto f(r+qx)$ gives, for $q>0$,
\begin{equation}\label{tools:poisson-progression}
 \sum_{n\in\mathbb Z}f(r+qn)
 =\frac1q\sum_{m\in\mathbb Z}
          e^{2\pi imr/q}\widehat f(m/q).
\end{equation}
Indeed, substitution in the Fourier integral gives its $m$th
coefficient as $q^{-1}e^{2\pi imr/q}\widehat f(m/q)$.

\subsection{The gamma function and its zeros}

For $\Re z>0$ define
\[
 \Gamma(z)=\int_0^\infty e^{-t}t^{z-1}\,dt.
\]
The integral and its derivatives converge locally uniformly in this
half-plane. Euler's constant is
$\gamma=\lim_{n\to\infty}(\sum_{k=1}^n k^{-1}-\log n)$.
The limit exists because the sequence is decreasing and positive, as
comparison with the integral of $1/x$ shows.

\begin{lemma}\label{tools:gamma}
The gamma function is nonzero for $\Re z>0$. Its reciprocal extends
to the entire function
\begin{equation}\label{tools:gamma-product}
 \frac1{\Gamma(z)}
   =ze^{\gamma z}\prod_{k=1}^{\infty}
        \left(1+\frac zk\right)e^{-z/k}.
\end{equation}
Its zeros are precisely $0,-1,-2,\ldots$, each simple.
Also $\Gamma(z+1)=z\Gamma(z)$ where both sides are defined, and
$\Gamma(z)\sim1/z$ as $z\to0$.
\end{lemma}
\begin{proof}
For $\Re z>0$, repeated integration by parts gives
\[
 \int_0^1 t^{z-1}(1-t)^n\,dt
  =\frac{n!}{z(z+1)\cdots(z+n)}.
\]
After multiplying by $n^z$ and putting $u=nt$, dominated convergence
applies because $(1-u/n)^n\leq e^{-u}$ for $0\leq u\leq n$.
It follows that
\[
 F_n(z):=\frac{n!n^z}{z(z+1)\cdots(z+n)}\longrightarrow\Gamma(z).
\]
On the other hand, the exact reciprocal is
\[
 F_n(z)^{-1}
  =z\exp\left(z\left(\sum_{k=1}^n\frac1k-\log n\right)\right)
       \prod_{k=1}^n(1+z/k)e^{-z/k}.
\]
On a compact set, the logarithms of the tail factors are uniformly
$O(k^{-2})$. The infinite product therefore converges locally uniformly;
off the nonpositive integers its tail and its finitely many initial
factors are nonzero. Thus the displayed reciprocals converge to a
nonzero function for $\Re z>0$. Passing to the limit in
$F_n(z)F_n(z)^{-1}=1$ proves both nonvanishing of $\Gamma$ and the
product formula there. The product itself is entire. At a nonpositive
integer exactly one factor vanishes simply, which proves the assertion
about its zeros and defines the meromorphic extension of $\Gamma$.
Finally, integration by parts in Euler's integral gives the recurrence.
Since $\Gamma(1)=1$, it also gives the asymptotic at zero.
\end{proof}

\section{Primitive characters and the critical strip}
\label{sec:primitive-reduction}

This section isolates a classical part of the argument which can be written
without any of the automorphic constructions in the announced proof.  It does
not establish the required nonvanishing in the right half-plane.  It proves,
with the analytic details made explicit, why that assertion would imply the
usual statement about zeros of every Dirichlet $L$-function in the critical
strip.

\subsection{Primitive Gauss sums}

Let $q$ be a positive integer.  A Dirichlet character modulo $q$ is a
homomorphism from $(\mathbb Z/q\mathbb Z)^\times$ to $\mathbb C^\times$,
extended periodically to the integers by zero on integers not relatively
prime to $q$.  For $q=1$ the unique character has value $1$ at every integer.
A character modulo $q$ is \emph{primitive} if it is not induced from a
character modulo a proper divisor of $q$.  Here induction means that the
values agree at all integers relatively prime to $q$.  The conductor of a
character is the modulus of its inducing primitive character.

Every character modulo $N$ is induced from a unique primitive character
modulo a divisor $q$ of $N$.  To see this explicitly, write
$N=\prod_{p\mid N}p^{e_p}$.  The Chinese remainder theorem identifies
its unit group with the product of the unit groups modulo $p^{e_p}$,
and the character is uniquely a product of characters of those factors.
For each $p$ choose the smallest $f_p\in\{0,\ldots,e_p\}$ through which
the corresponding character factors under reduction modulo $p^{f_p}$;
the unit group for $p^0=1$ is trivial.  Reduction maps of these unit groups
are surjective, so the factored characters are unique.  Their product
has modulus $q=\prod p^{f_p}$, induces the original character, and is
primitive by minimality at each prime.  Any primitive inducing modulus
must have exactly these exponents: an exponent below $f_p$ is impossible,
and one above $f_p$ gives an imprimitive character.  This proves the
assertion and the uniqueness.

Every character value at a unit is a root of unity, because the group of
units is finite.  Thus $|\psi(n)|\leq 1$ and
$\overline{\psi(n)}=\psi(n)^{-1}$ for units $n$.
For a primitive character $\psi$ modulo $q>1$, put
\[
 a=\begin{cases}0,&\psi(-1)=1,\\1,&\psi(-1)=-1,\end{cases}
 \qquad
 \tau(\psi)=\sum_{r\bmod q}\psi(r)e^{2\pi i r/q}.
\]
The two possibilities for $\psi(-1)$ follow from $(-1)^2=1$.

\begin{lemma}\label{prim:gauss}
For every integer $m$,
\[
 \sum_{r\bmod q}\psi(r)e^{2\pi i mr/q}
   =\overline{\psi(m)}\tau(\psi),
 \qquad |\tau(\psi)|=\sqrt q.
\]
\end{lemma}

\begin{proof}
If $m$ is relatively prime to $q$, multiplication by $m$ permutes the
residue classes modulo $q$.  Substituting $u=mr$ gives the first identity.

Suppose now that $d=(m,q)>1$, and put $q'=q/d$.  Reduction of units modulo
$q$ onto the units modulo $q'$ is surjective.  Indeed, a residue class
coprime to $q'$ can be lifted to an integer coprime also to every prime
dividing $q$ but not $q'$, by the Chinese remainder theorem.  If $\psi$
were trivial on the kernel of this reduction, it would factor through
the units modulo $q'$, contradicting primitivity.  There is therefore a
unit $u$ modulo $q$ such that $u\equiv 1\pmod {q'}$ and $\psi(u)\neq1$.
Since $m(u-1)$ is divisible by $q$, substitution $r\mapsto ur$ in the sum
shows that the sum equals $\psi(u)$ times itself.  It is zero.  The right
side is also zero because $\psi(m)=0$.

Write $S(m)$ for the sum in the first identity.  Finite Fourier
orthogonality gives
\begin{align*}
 \sum_{m\bmod q}|S(m)|^2
 &=\sum_{r,t\bmod q}\psi(r)\overline{\psi(t)}
       \sum_{m\bmod q}e^{2\pi i m(r-t)/q}\\
 &=q\sum_{r\bmod q}|\psi(r)|^2=q\varphi(q),
\end{align*}
where $\varphi(q)$ is the number of units modulo $q$.
On the other hand, the identity already proved makes the left side
$|\tau(\psi)|^2\varphi(q)$.  Division by the positive integer
$\varphi(q)$ proves the second assertion.
\end{proof}

\subsection{The primitive functional equation}

For a positive real number $x$, complex powers mean
$x^s=\exp(s\log x)$ with the real logarithm.
For $\operatorname{Re}s>1$, define
\[
 L(s,\psi)=\sum_{n=1}^{\infty}\frac{\psi(n)}{n^s}.
\]
The series converges absolutely and locally uniformly in that half-plane.
For primitive $\psi$ modulo $q>1$, define its completion there by
\[
 \Lambda(s,\psi)
 =\left(\frac q\pi\right)^{(s+a)/2}
     \Gamma\left(\frac{s+a}{2}\right)L(s,\psi).
\]
Here $\Gamma$ is Euler's gamma function. Lemma~\ref{tools:gamma}
proves its nonvanishing and the entire continuation of its reciprocal.
We use Poisson summation in the normalization of
Lemma~\ref{tools:poisson} and equation~\eqref{tools:poisson-progression}.

\begin{proposition}\label{prim:functional-equation}
If $\psi$ is primitive modulo $q>1$, then $\Lambda(s,\psi)$ extends to an
entire function and
\[
 \Lambda(s,\psi)
    =\varepsilon_\psi\Lambda(1-s,\overline\psi),
 \qquad
 \varepsilon_\psi=i^{-a}\frac{\tau(\psi)}{\sqrt q},
 \qquad |\varepsilon_\psi|=1.
\]
The unique primitive principal character has conductor $1$.  Its
$L$-function is $\zeta(s)$, and
\[
 \Lambda_0(s)=\pi^{-s/2}\Gamma(s/2)\zeta(s)
\]
has a meromorphic continuation with only simple poles at $0$ and $1$,
and satisfies $\Lambda_0(s)=\Lambda_0(1-s)$.
\end{proposition}

\begin{proof}
First assume $q>1$ and put
\[
 \theta_\psi(t)=\sum_{n\in\mathbb Z}
      \psi(n)n^a e^{-\pi n^2t/q},\qquad t>0.
\]
When $a=0$ the factor $n^a$ is interpreted as $1$; the term $n=0$ is
still zero because $\psi(0)=0$.  The series and all its derivatives
converge uniformly on compact subintervals of $(0,\infty)$.

Applying Poisson summation in each residue class modulo $q$ yields
\[
 \sum_{n\in\mathbb Z}\psi(n)f(n)
   =\frac1q\sum_{m\in\mathbb Z}\widehat f(m/q)
       \sum_{r\bmod q}\psi(r)e^{2\pi i mr/q}
   =\frac{\tau(\psi)}q\sum_{m\in\mathbb Z}
       \overline{\psi(m)}\widehat f(m/q).
\]
Use $f(x)=x^a e^{-\pi tx^2/q}$.  The Gaussian Fourier transform and its
first derivative give, for $a=0,1$,
\[
 \widehat f(y)=(-i)^a
      (q/t)^{a+1/2}y^a e^{-\pi qy^2/t}.
\]
Consequently
\begin{equation}\label{prim:theta-transform}
 \theta_\psi(t)
  =\varepsilon_\psi t^{-a-1/2}
       \theta_{\overline\psi}(1/t).
\end{equation}
Lemma~\ref{prim:gauss} proves $|\varepsilon_\psi|=1$.

The theta series is exponentially small as $t\to\infty$.
Equation~\eqref{prim:theta-transform} shows that as $t\to0^+$ it is
bounded by a fixed power of $t^{-1}$ times $e^{-c/t}$, for some $c>0$.
It follows that
\[
 F_\psi(s)=\frac12\int_0^\infty
           \theta_\psi(t)t^{(s+a)/2-1}\,dt
\]
converges absolutely and locally uniformly for every $s\in\mathbb C$.
The same is true after inserting any fixed power of $\log t$, so
differentiation under the integral proves that $F_\psi$ is entire.

For $\operatorname{Re}s>1$, absolute convergence justifies termwise
integration.  Indeed, after taking absolute values, the integral of the
$n$th term is a constant depending on $s$ times $|n|^{-\operatorname{Re}s}$.
Since $\psi(-n)(-n)^a=\psi(n)n^a$, the factor $1/2$ removes the duplicate
negative terms.  Substitution $u=\pi n^2t/q$ therefore gives
\[
 F_\psi(s)
 =\Gamma\left(\frac{s+a}{2}\right)
      \left(\frac q\pi\right)^{(s+a)/2}
      \sum_{n\geq1}\frac{\psi(n)}{n^s}
 =\Lambda(s,\psi).
\]
Finally insert~\eqref{prim:theta-transform} in the integral and make the
change of variable $u=1/t$.  Absolute convergence at both endpoints
justifies this substitution and yields
\[
 F_\psi(s)=\frac{\varepsilon_\psi}{2}
       \int_0^\infty\theta_{\overline\psi}(u)
           u^{(1-s+a)/2-1}\,du
   =\varepsilon_\psi F_{\overline\psi}(1-s).
\]
This proves the functional equation and the continuation for $q>1$.
Dividing the entire completion by the displayed gamma and power factors
also continues $L(s,\psi)$; equivalently, multiplication by the entire
function $1/\Gamma((s+a)/2)$ shows that $L(s,\psi)$ is entire.

For the principal character of conductor $1$, let
$\theta(t)=\sum_{n\in\mathbb Z}e^{-\pi n^2t}$ and $A(t)=\theta(t)-1$.
Poisson summation gives $\theta(t)=t^{-1/2}\theta(1/t)$.
For $\operatorname{Re}s>1$, termwise integration gives
$\Lambda_0(s)=\frac12\int_0^\infty A(t)t^{s/2-1}\,dt$.
Split the integral at $1$ and transform the portion from $0$ to $1$.
The result is
\begin{equation}\label{prim:zeta-continuation}
 \Lambda_0(s)=\frac1{s-1}-\frac1s
  +\frac12\int_1^\infty A(t)
       \bigl(t^{s/2-1}+t^{(1-s)/2-1}\bigr)\,dt.
\end{equation}
The integral is entire by exponential decay, and the right side is
unchanged by $s\mapsto1-s$.  It has exactly the asserted poles, with
residues $1$ at $1$ and $-1$ at $0$.
\end{proof}

\subsection{Changing the modulus}

Suppose $\chi$ modulo $N$ is induced by a primitive character $\psi$
modulo $q$, where $q\mid N$.  Define the entire function
\[
 E_{N,q}(s,\psi)
   =\prod_{\substack{p\mid N\\p\nmid q}}
         (1-\psi(p)p^{-s}),
\]
where $p$ runs over primes and an empty product is $1$.

\begin{lemma}\label{prim:change-level}
As an identity of meromorphic functions,
\[
 L(s,\chi)=L(s,\psi)E_{N,q}(s,\psi).
\]
Moreover, $E_{N,q}(s,\psi)\neq0$ whenever
$\operatorname{Re}s>0$.  Thus $L(s,\chi)$ and $L(s,\psi)$ have the same
zeros, with the same multiplicities, in that open half-plane, away from
any pole.
\end{lemma}

\begin{proof}
In $\operatorname{Re}s>1$, absolute convergence and unique prime
factorization give the Euler products
\[
 L(s,\chi)=\prod_p(1-\chi(p)p^{-s})^{-1},
 \qquad
 L(s,\psi)=\prod_p(1-\psi(p)p^{-s})^{-1}.
\]
For primes $p\nmid N$ the character values agree.  For primes dividing
$q$ both values vanish.  For $p\mid N$, $p\nmid q$, the factor of
$L(s,\chi)$ is $1$, whereas the factor of $L(s,\psi)$ is
$(1-\psi(p)p^{-s})^{-1}$.  Cancelling these factors proves the identity
in that half-plane.  The identity supplies the meromorphic continuation
of $L(s,\chi)$ and remains valid everywhere by uniqueness of analytic
continuation.

If $\sigma=\operatorname{Re}s>0$, then
\[
 |\psi(p)p^{-s}|\leq p^{-\sigma}<1,
 \qquad
 |1-\psi(p)p^{-s}|\geq1-p^{-\sigma}>0.
\]
Every factor, and hence the finite product, is nonzero.  Multiplication
by a holomorphic function which is nonzero at a point preserves the
order of a zero there, proving the multiplicity statement.
\end{proof}

For the primitive principal character, we use the conventions
$a=0$ and $\Lambda(s,\psi)=\Lambda_0(s)$ from
Proposition~\ref{prim:functional-equation}.

The restriction $\operatorname{Re}s>0$ in this lemma is essential.
For example, for the principal character modulo a prime $p$,
$L(s,\chi)=\zeta(s)(1-p^{-s})$ as meromorphic functions.
The extra factor vanishes at $s=2\pi i k/\log p$ for $k\in\mathbb Z$.
In particular, $s=0$ is an additional zero since $\zeta(0)=-1/2$.
This value follows directly from~\eqref{prim:zeta-continuation}:
$\Lambda_0(s)\sim-1/s$ and $\Gamma(s/2)\sim2/s$ at zero.
Thus imprimitive characters require care on the boundary of the critical
strip.

\subsection{The precise reduction}

For clarity, state the needed half-plane assertion explicitly:
\begin{quote}
For every positive integer $N$, every Dirichlet character $\chi$ modulo
$N$, and every $s\in\mathbb C$ with $\operatorname{Re}s>1/2$, the value
$L(s,\chi)$ is nonzero whenever $\chi$ is nonprincipal or $s\neq1$.
\end{quote}
At $s=1$ a principal $L$-function has a pole, so that exception is part of
the statement, not a possible zero.  The value of a meromorphic function
at its pole is not used in the argument.

\begin{theorem}[Critical-strip reduction]\label{prim:strip-reduction}
Assume the half-plane assertion above.  Then for every Dirichlet
character $\chi$ modulo a positive integer $N$,
\[
 0<\operatorname{Re}s<1,\qquad L(s,\chi)=0
 \quad\Longrightarrow\quad
 \operatorname{Re}s=\frac12.
\]
In fact, it suffices to assume nonvanishing for primitive characters at
points satisfying $1/2<\operatorname{Re}s<1$.
\end{theorem}

\begin{proof}
Let $\psi$ be the primitive character inducing $\chi$.
Lemma~\ref{prim:change-level} shows that a zero $s$ of $L(s,\chi)$ in
the strip is a zero of $L(s,\psi)$.  If
$\operatorname{Re}s>1/2$, the stated nonvanishing already gives a
contradiction.

Suppose instead that $\operatorname{Re}s<1/2$.  The gamma factor and
the power of $q/\pi$ in the completion are finite and nonzero at $s$,
because $\operatorname{Re}((s+a)/2)>0$.  Therefore
$\Lambda(s,\psi)=0$.  Proposition~\ref{prim:functional-equation} and
$\varepsilon_\psi\neq0$ imply
$\Lambda(1-s,\overline\psi)=0$.  For the character of conductor $1$
the same inference uses~\eqref{prim:zeta-continuation}; neither $s$ nor
$1-s$ is a pole.  The conjugate character is primitive: any factorization
of $\overline\psi$ through a proper divisor, when conjugated, would give
one for $\psi$.  Its parity is the same as that of $\psi$.

At $1-s$ the gamma and power factors are again finite and nonzero,
since $0<\operatorname{Re}(1-s)<1$.  Hence
$L(1-s,\overline\psi)=0$.  But
$1/2<\operatorname{Re}(1-s)<1$, contradicting the assumed primitive
nonvanishing.  The only remaining possibility is
$\operatorname{Re}s=1/2$.
\end{proof}

The character changes to its conjugate in this reflection.  For a
nonreal character it is generally incorrect to replace
$\overline\psi$ by $\psi$ in the displayed functional equation.
No bound on the conductor, and no restriction to real or even
characters, is present in the theorem.

The reduction just proved is the classical part of the passage from
half-plane nonvanishing to the conclusion of the announced
result~\cite{AlpogeSource}. The half-plane estimate itself requires
further arithmetic input.

\section{Smoothed prime sums and nonvanishing}
\label{sec:prime-sum-reduction}

The purpose of this section is to identify a precise analytic estimate
which suffices for nonvanishing.  A global upper bound for a certain real
combination of smoothed prime sums is sufficient; a corresponding lower
bound is not required.  This is a conditional reduction.  No such upper
bound is established here.

\subsection{Normalization and Mellin transforms}

Let $\Lambda(n)$ denote the von Mangoldt function: it is $\log p$ when
$n=p^k$ for a prime $p$ and an integer $k\geq1$, and is zero otherwise.
Let $\gamma=\lim_{M\to\infty}(\sum_{n=1}^M n^{-1}-\log M)$.
For $x\geq1$, define
\begin{align}
 S(x)&=\sum_{n\leq x}\Lambda(n)(x/n-1),\label{mellin:S}\\
 G(x)&=S(x)-x\log x+(1+\gamma)x.\label{mellin:G}
\end{align}
For a nonprincipal Dirichlet character $\chi$, assume explicitly that
$L(1,\chi)\neq0$, and set
\begin{align}
 c_\chi&=\frac{L'(1,\chi)}{L(1,\chi)},\label{mellin:cchi}\\
 g_\chi(x)&=\sum_{n\leq x}\Lambda(n)\chi(n)(x/n-1)+c_\chi x,
 \qquad G_\chi(x)=\operatorname{Re}g_\chi(x).\label{mellin:gchi}
\end{align}
The assumption at $1$ is the classical nonvanishing theorem for
nonprincipal Dirichlet characters.  In the conditional statements below
it is an explicit hypothesis.  Each function in
\eqref{mellin:S}--\eqref{mellin:gchi} is continuous on $[1,\infty)$:
the newly appearing term at an integer $n$ has value zero at that point.
The elementary inequalities $|\chi(n)|\leq1$ and
$0\leq\Lambda(n)\leq\log n$ give
\begin{equation}\label{mellin:trivial-growth}
 |G(x)|+|g_\chi(x)|=O_\chi\bigl(x\log^2(2x)\bigr).
\end{equation}
For example, the uncorrected sum is bounded in absolute value by
$x\sum_{n\leq x}(\log n)/n$.

We record explicitly the analytic information needed near $s=1$.
For $\operatorname{Re}s>1$, partial summation gives
\begin{equation}\label{mellin:zeta-fractional}
 \zeta(s)=\frac{s}{s-1}
       -s\int_1^\infty\{x\}x^{-s-1}\,dx,
\end{equation}
where $\{x\}=x-\lfloor x\rfloor$.  Since $0\leq\{x\}<1$, the
integral is holomorphic for $\operatorname{Re}s>0$, by locally uniform
absolute convergence, also after differentiation.  This proves a
meromorphic continuation of $\zeta$ to that half-plane with just a simple
pole at $1$.  At an integer $M\geq1$,
\[
 \int_1^M\frac{\{x\}}{x^2}\,dx
   =\log M-\sum_{n=1}^{M-1}n\left(\frac1n-\frac1{n+1}\right)
   =1+\log M-\sum_{n=1}^M\frac1n.
\]
Consequently the constant term of~\eqref{mellin:zeta-fractional} at
$1$ is $\gamma$, and
\begin{equation}\label{mellin:zeta-laurent}
 \zeta(s)=\frac1{s-1}+\gamma+O(s-1),\qquad
 Z(s):=(s-1)\zeta(s),\quad Z(1)=1,\quad Z'(1)=\gamma.
\end{equation}
The harmonic limit defining $\gamma$ exists by integral comparison;
$\sum_{n=1}^M n^{-1}-\log M$ is positive and decreasing.

For a nonprincipal character, the sum of its values in a complete period
is zero.  Indeed, multiplication by a unit on which the character is
not $1$ permutes that sum and multiplies it by a different scalar.
Thus $A_\chi(x)=\sum_{n\leq x}\chi(n)$ is bounded.  The formula
\begin{equation}\label{mellin:dirichlet-continuation}
 L(s,\chi)=s\int_1^\infty A_\chi(x)x^{-s-1}\,dx
\end{equation}
continues $L(s,\chi)$ holomorphically to $\operatorname{Re}s>0$.
The equality first follows by summation by parts for real part greater
than $1$, and the integral converges locally uniformly for real part
greater than $0$.  These elementary continuations agree with the usual
$L$-functions by uniqueness of analytic continuation.

\begin{lemma}\label{mellin:identities}
For $\operatorname{Re}s>1$ the following integrals converge absolutely:
\begin{align}
 \int_1^\infty G(x)x^{-s-1}\,dx
 &=\frac{-\zeta'(s)/\zeta(s)}{s(s-1)}
     -\frac1{(s-1)^2}+\frac{1+\gamma}{s-1},\label{mellin:G-transform}\\
 \int_1^\infty g_\chi(x)x^{-s-1}\,dx
 &=\frac{-L'(s,\chi)/L(s,\chi)}{s(s-1)}
     +\frac{c_\chi}{s-1}.\label{mellin:g-transform}
\end{align}
The expression in~\eqref{mellin:G-transform} has a removable singularity
at $1$.  The same is true of~\eqref{mellin:g-transform} when
$L(1,\chi)\neq0$.
\end{lemma}

\begin{proof}
For $n\geq1$ and $\operatorname{Re}s>1$, direct integration gives
\[
 \int_n^\infty(x/n-1)x^{-s-1}\,dx
   =\frac{n^{-s}}{s(s-1)}.
\]
The sum of the absolute values of the corresponding integrals, with
coefficients bounded by $\Lambda(n)$, is finite.  This follows from
$\sum_{n\geq2}(\log n)n^{-\operatorname{Re}s}<\infty$.
Fubini's theorem therefore permits the exchange of the sum and the
integral.  The absolutely convergent Euler products give
\[
 -\frac{\zeta'(s)}{\zeta(s)}
   =\sum_{n\geq1}\frac{\Lambda(n)}{n^s},\qquad
 -\frac{L'(s,\chi)}{L(s,\chi)}
   =\sum_{n\geq1}\frac{\Lambda(n)\chi(n)}{n^s}.
\]
One obtains these identities by expanding the logarithm of each Euler
factor and differentiating its locally uniformly convergent series.
The same absolute convergence proves that the Euler products are
nonzero in this half-plane.
Finally,
\[
 \int_1^\infty x\,x^{-s-1}\,dx=\frac1{s-1},\qquad
 \int_1^\infty x\log x\,x^{-s-1}\,dx=\frac1{(s-1)^2}.
\]
These calculations prove both identities.

For the assertion at $1$, write $u=s-1$.  Equation
\eqref{mellin:zeta-laurent} implies
$-\zeta'/\zeta=1/u-\gamma+O(u)$.  Hence
\[
 \frac{-\zeta'/\zeta}{s(s-1)}
   =\frac1{u^2}-\frac{1+\gamma}{u}+O(1),
\]
so the two displayed correction terms cancel both polar terms.
For $\chi$, the quotient $L'/L$ is analytic near $1$ with value
$c_\chi$.  The polar term in its negative divided by $s(s-1)$ is
$-c_\chi/(s-1)$, again exactly cancelled by the correction term.
\end{proof}

The coefficient $1+\gamma$ and the sign of $c_\chi$ in these definitions
are therefore forced by cancellation at $1$.  Altering either linear
term changes the problem being estimated.

\subsection{The real singularity of a nonnegative Mellin transform}

The following version of Landau's principle supplies the analytic
reason that a one-sided bound can be enough.

\begin{lemma}\label{mellin:landau}
Let $H:[1,\infty)\to[0,\infty)$ be measurable and locally integrable,
and suppose $\int_1^\infty H(x)x^{-\sigma-1}\,dx$ is finite for at least
one real $\sigma$.  Put
\[
 \sigma_c=\inf\left\{\sigma\in\mathbb R:
       \int_1^\infty H(x)x^{-\sigma-1}\,dx<\infty\right\}.
\]
If $\sigma_c$ is finite, the holomorphic function
\[
 F(s)=\int_1^\infty H(x)x^{-s-1}\,dx,
       \qquad \operatorname{Re}s>\sigma_c,
\]
has no holomorphic continuation to a neighborhood of the real point
$\sigma_c$.
\end{lemma}

\begin{proof}
Nonnegativity makes the set of real exponents for which the integral
is finite an upper interval.  For every $\sigma>\sigma_c$, one can
choose a smaller exponent in that interval; hence the integral is
finite at $\sigma$.  Local uniform domination then makes $F$ holomorphic
in the indicated half-plane.  For a real $a>\sigma_c$,
\begin{equation}\label{mellin:moments}
 F^{(k)}(a)=(-1)^k\int_1^\infty
          H(x)(\log x)^k x^{-a-1}\,dx.
\end{equation}
Differentiation is justified by comparison with a slightly smaller
exponent greater than $\sigma_c$: every fixed power of $\log x$ is
bounded by a constant times an arbitrarily small positive power of $x$.

Suppose now that $F$ continues holomorphically to a disk
$|s-\sigma_c|<r$, for some $r>0$.  Set
$a=\sigma_c+r/4$ and $h=r/2$.  The Taylor series of this continuation
about $a$ converges at $a-h$, since the disk about $a$ of radius $3r/4$
is contained in the first disk.  Using~\eqref{mellin:moments}, all the
terms in this Taylor expansion are nonnegative real numbers:
\begin{align*}
 F(a-h)
  &=\sum_{k=0}^{\infty}\frac{(-h)^kF^{(k)}(a)}{k!}\\
  &=\sum_{k=0}^{\infty}\frac{h^k}{k!}
       \int_1^\infty H(x)(\log x)^k x^{-a-1}\,dx\\
  &=\int_1^\infty H(x)x^{-a-1}
       \sum_{k=0}^{\infty}\frac{(h\log x)^k}{k!}\,dx\\
  &=\int_1^\infty H(x)x^{-(a-h)-1}\,dx.
\end{align*}
The exchange uses Tonelli's theorem, so no integrability assumption
below $\sigma_c$ has been inserted.  The finite Taylor sum limit proves
that the last integral is finite.  But
$a-h=\sigma_c-r/4<\sigma_c$, contradicting the definition of
$\sigma_c$.
\end{proof}

\subsection{A global one-sided criterion for a family}

Let $\mathcal X$ be a finite set of nonprincipal Dirichlet characters,
possibly of different moduli.  Suppose that $\mathcal X$ is closed
under complex conjugation, and assign positive integer multiplicities
$m_\chi$ satisfying $m_{\overline\chi}=m_\chi$.  Assume
$L(1,\chi)\neq0$ for every $\chi\in\mathcal X$.  Let $R$ be a positive
integer and define the real-valued function
\begin{equation}\label{mellin:B}
 B(x)=R G(x)+\sum_{\chi\in\mathcal X}m_\chi G_\chi(x).
\end{equation}
The empty family is allowed.  It yields a criterion for $\zeta$ alone.

\begin{theorem}\label{mellin:one-sided-family}
Let $0\leq\theta<1$.  If there are constants $C_0\in\mathbb R$ and
$x_0\geq1$ such that
\begin{equation}\label{mellin:global-bound}
 B(x)\leq C_0x^\theta\qquad\text{for every }x\geq x_0,
\end{equation}
then $\zeta$ has no zeros in $\operatorname{Re}s>\theta$, and
$L(s,\chi)\neq0$ in this half-plane for every
$\chi\in\mathcal X$.  The pole of $\zeta$ at $1$ is understood as a
pole, not as a value to which the assertion is applied.
\end{theorem}

\begin{proof}
The identity
$L(s,\overline\chi)=\overline{L(\overline s,\chi)}$ follows first
from the Dirichlet series and then by analytic continuation.  In
particular $c_{\overline\chi}=\overline{c_\chi}$.  Therefore
\[
 \sum_{\chi\in\mathcal X}m_\chi G_\chi(x)
   =\sum_{\chi\in\mathcal X}m_\chi g_\chi(x),
 \qquad
 c=\sum_{\chi\in\mathcal X}m_\chi c_\chi\in\mathbb R.
\]
In the half-plane $\operatorname{Re}s>0$, form the holomorphic function
\[
 P(s)=Z(s)^R\prod_{\chi\in\mathcal X}L(s,\chi)^{m_\chi}.
\]
It is not identically zero, since the Euler products are nonzero for
$\operatorname{Re}s>1$.  Also $P(1)\neq0$ by the assumptions at $1$.
Put
\[
 a=-R,\qquad b=R(1+\gamma)+c.
\]
Logarithmic differentiation at $1$ gives
\begin{equation}\label{mellin:P-at-one}
 \frac{P'(1)}{P(1)}=R\gamma+c=a+b.
\end{equation}
Adding the identities in Lemma~\ref{mellin:identities} and using
$Z'/Z=1/(s-1)+\zeta'/\zeta$ yields, for
$\operatorname{Re}s>1$,
\begin{equation}\label{mellin:B-transform}
 \int_1^\infty B(x)x^{-s-1}\,dx
   =\frac{-P'(s)/P(s)+a+bs}{s(s-1)}.
\end{equation}

By continuity on $[1,x_0]$, increase $C_0$ to a nonnegative constant
$C$ so that
$H(x)=Cx^\theta-B(x)\geq0$ for all $x\geq1$.
Equation~\eqref{mellin:trivial-growth} shows that the Mellin transform
of $H$ converges absolutely for real part greater than $1$.
Let $\sigma_c\leq1$ be its abscissa of convergence; the case
$\sigma_c=-\infty$ is permitted.  For real part greater than $1$,
\begin{equation}\label{mellin:H-transform}
 F(s):=\int_1^\infty H(x)x^{-s-1}\,dx
   =\frac{P'(s)/P(s)-(a+bs)}{s(s-1)}+\frac C{s-\theta}.
\end{equation}
Denote the meromorphic expression on the right by $\Phi(s)$.
It is meromorphic for $\operatorname{Re}s>0$.
Equation~\eqref{mellin:P-at-one} shows that its singularity at $1$ is
removable; here $\theta<1$ and $P(1)\neq0$ are used.

We claim that $\sigma_c\leq\theta$.  Otherwise
$\theta<\sigma_c\leq1$ is a finite real number, and the identity theorem
for meromorphic functions gives $F=\Phi$ throughout
$\operatorname{Re}s>\sigma_c$.  If $\sigma_c=1$, the removability just
proved contradicts Lemma~\ref{mellin:landau}.
Suppose therefore that $\beta=\sigma_c\in(\theta,1)$.
If $P(\beta)\neq0$, then $\Phi$ is holomorphic near $\beta$, again
contradicting that lemma.
If $P(\beta)=0$, let $m\geq1$ be its order of vanishing.  Locally,
$P(s)=(s-\beta)^mU(s)$ with $U(\beta)\neq0$, so
\[
 \frac{P'(s)}{P(s)}=\frac m{s-\beta}+\frac{U'(s)}{U(s)}.
\]
Thus the residue of $\Phi$ at $\beta$ is
\[
 \frac{m}{\beta(\beta-1)}<0.
\]
As real $\sigma$ decreases to $\beta$ from above,
$F(\sigma)=\Phi(\sigma)\to-\infty$.  This is impossible because
$F(\sigma)$ is the integral of a nonnegative real function.
The claim follows.

Consequently $F$ is holomorphic in $\operatorname{Re}s>\theta$.
Define there the holomorphic coefficient
\[
 Q(s)=a+bs+s(s-1)\left(F(s)-\frac C{s-\theta}\right).
\]
Multiplication of~\eqref{mellin:H-transform} by $s(s-1)P(s)$ gives
$P'=QP$ for real part greater than $1$.  Both sides are holomorphic
in the connected half-plane $\operatorname{Re}s>\theta$, so the
identity theorem extends this equality throughout that half-plane.
A nonzero holomorphic function satisfying $P'=QP$ cannot have a zero:
at a zero of order $m\geq1$, the order of $P'$ would be $m-1$, whereas
the order of $QP$ would be at least $m$.  Hence $P$ has no zeros there.
Every exponent in its product is a positive integer.  It follows that
each $L(s,\chi)$ and $Z(s)$ is nonzero there.  Finally
$Z(s)=(s-1)\zeta(s)$ away from $1$, proving the assertion for $\zeta$.
\end{proof}

The argument uses no estimate for $-B(x)$.
Its essential sign is the negative residue
$m/(\beta(\beta-1))$ for $0<\beta<1$.  Positive multiplicities prevent
cancellation of zero orders in the product $P$.

\subsection{An absolute-bound criterion and the limitation of windows}

\begin{corollary}\label{mellin:absolute}
Fix $\theta\geq0$.  If, for every $\epsilon>0$,
$G(x)=O_\epsilon(x^{\theta+\epsilon})$ in absolute value, then
$\zeta$ has no zeros with real part greater than $\theta$.
For a fixed nonprincipal character with $L(1,\chi)\neq0$, if
$g_\chi(x)=O_\epsilon(x^{\theta+\epsilon})$ in absolute value for
every $\epsilon>0$, then $L(s,\chi)$ has no zeros in that half-plane.
\end{corollary}

\begin{proof}
The assumptions make the respective Mellin integrals $I(s)$ and $J(s)$
holomorphic in $\operatorname{Re}s>\theta$.  Indeed, on any compact
subset choose $\epsilon$ smaller than its distance from the boundary;
the resulting power bound dominates the integrands and all their
derivatives locally uniformly.  The identities of
Lemma~\ref{mellin:identities}, multiplied out, become
\[
 Z'(s)=\bigl[-1+(1+\gamma)s-s(s-1)I(s)\bigr]Z(s),
\]
and
\[
 L'(s,\chi)=\bigl[c_\chi s-s(s-1)J(s)\bigr]L(s,\chi).
\]
They initially hold for real part greater than $\max(1,\theta)$ and
extend to $\operatorname{Re}s>\theta$ by the identity theorem.
The same order-of-vanishing argument as above excludes zeros.
The elementary continuations are holomorphic throughout this
half-plane because $\theta\geq0$.
\end{proof}

The global quantifier in~\eqref{mellin:global-bound} must be preserved.
An estimate only on arbitrarily large intervals
$[Y,Y^\lambda]$, with fixed $\lambda>1$, does not by itself make
$Cx^\theta-B(x)$ nonnegative on the entire integration range.  It
therefore cannot be substituted into the proof of
Theorem~\ref{mellin:one-sided-family}.

For a general continuous function, even an absolute bound on such
intervals gives no global growth estimate.  To see this, choose $Y_1>1$ and numbers
$Y_j\to\infty$ so rapidly that
$Y_{j+1}>4Y_j^\lambda$.  Let $\eta$ be a continuous nonnegative function
supported on $[2,3]$, with $\eta(5/2)=1$, and fix $\beta>\theta$.
The locally finite sum
\[
 f(x)=\sum_{j\geq1}x^\beta\eta(x/Y_j^\lambda)
\]
vanishes on every $[Y_j,Y_j^\lambda]$, while
$f((5/2)Y_j^\lambda)=((5/2)Y_j^\lambda)^\beta$.
It is not $O(x^\theta)$.  Moreover, if $\eta$ is positive on a
subinterval of $(2,3)$, its Mellin integral diverges for real exponents
at most $\beta$: each separate bump contributes a positive constant
times $(Y_j^\lambda)^{\beta-\sigma}$.

This example is not a counterexample to a theorem about the particular
prime sums above.  Their logarithmic-derivative structure may allow a
stronger theorem from long multiplicative intervals.  Such a theorem
requires an additional argument.  The global criterion proved here
does not establish it, and no global estimate is inferred from a
window estimate in this section.

% Independently written reductions for integration into the GRH paper.
% Required packages: amsmath, amssymb, amsthm.
% Required theorem environments: theorem, lemma, proposition, corollary,
% definition, remark. Bibliography keys: OpenAIOctober, OpenAISeptember.

\section{Smoothed sums and zeros of L-functions}
\label{qr-sec-mellin}

The final analytic step can be separated from the arithmetic estimates.
This is useful because a proposed simplification must supply a precisely
stated estimate; an analogy with a known proof does not supply that estimate.
We first prove the implication that we shall need. The argument is related
to the reduction in \cite[Section~3]{OpenAIOctober}.
% Source labels: sec:reduction, eq:mobius-saving.

Let $K$ be a number field, let $\mathcal O_K$ be its ring of integers,
and let $N\mathfrak a=|\mathcal O_K/\mathfrak a|$ denote the norm of a
nonzero integral ideal. Here a number field is a finite extension of $\mathbb Q$, and
$\mathcal O_K$ consists of its elements integral over $\mathbb Z$.
For clarity, a ray class group is the group of fractional ideals
coprime to a fixed nonzero ideal $\mathfrak f$, modulo the principal
ideals $(\alpha)$ for which $\alpha\equiv1\pmod{\mathfrak f}$
locally at each prime dividing $\mathfrak f$, with positivity imposed
at a chosen set of real embeddings. Locally the congruence means
$v_{\mathfrak p}(\alpha-1)\geq v_{\mathfrak p}(\mathfrak f)$,
where $v_{\mathfrak p}$ is the exponent of $\mathfrak p$ in a
principal fractional ideal.
A finite-order Hecke character $\nu$ is a homomorphism of such a
ray class group into the complex roots of unity. We view it as a
completely multiplicative function on integral ideals coprime to
$\mathfrak f$ and extend it by zero to all other ideals. The least
possible finite modulus is its conductor, and we use that modulus.
In particular, $|\nu(\mathfrak a)|\leq1$. The ideal M\"obius function is
defined by $\mu_K(\mathcal O_K)=1$, by $\mu_K(\mathfrak a)=0$ if a square
of a prime ideal divides $\mathfrak a$, and by
$\mu_K(\mathfrak p_1\cdots\mathfrak p_r)=(-1)^r$ for distinct prime ideals.
We use the classical analytic properties of the associated $L$-function:
its Euler product converges absolutely in $\Re s>1$, and it has a
meromorphic continuation. The result below is an implication from these
properties and a stated sum estimate; it does not prove that estimate.

For a finite set $S$ of prime ideals, write $(\mathfrak a,S)=1$ if none
of them divides $\mathfrak a$, and put
\begin{equation}\label{qr-eq-deleted-L}
 L_K^S(s,\nu)=L_K(s,\nu)
   \prod_{\mathfrak p\in S}
       \bigl(1-\nu(\mathfrak p)(N\mathfrak p)^{-s}\bigr).
\end{equation}
Every factor in this finite product is nonzero in $\Re s>0$.
For $W\in C_c^\infty((0,\infty);\mathbb C)$ and $X>0$, define
\begin{equation}\label{qr-eq-smooth-sum}
 A_W(X)=\sum_{(\mathfrak a,S)=1}
       \mu_K(\mathfrak a)\nu(\mathfrak a)
       W(N\mathfrak a/X).
\end{equation}
The sum is finite, since the norm is bounded on its support.

\begin{theorem}[A smoothed cancellation criterion]\label{qr-thm-mellin}
Let $\theta\geq0$. Suppose that for every $\varepsilon>0$ and every
$W\in C_c^\infty((0,\infty);\mathbb C)$ there is a constant
$C_{\nu,S,W,\varepsilon}$ such that
\begin{equation}\label{qr-eq-cancellation-input}
 |A_W(X)|\leq C_{\nu,S,W,\varepsilon}X^{\theta+\varepsilon}
 \qquad(X\geq2).
\end{equation}
Then $L_K(s,\nu)$ has no zero in $\Re s>\theta$.
\end{theorem}

\begin{proof}
Fix a test $W$ supported in $[a,b]\subset(0,\infty)$. Since a nonzero
integral ideal has norm at least one, $A_W(X)=0$ for $X<1/b$. Define
\[
 \mathcal M_W(s)=\int_0^\infty A_W(X)X^{-s}\,\frac{dX}{X}.
\]
This integral is locally uniformly convergent in $\Re s>\theta$.
Indeed, for a compact subset of that half-plane choose $\delta>0$ such
that $\Re s\geq\theta+2\delta$ there, and apply
\eqref{qr-eq-cancellation-input} with $\varepsilon=\delta$.
The integrand at infinity is then bounded by a constant times
$X^{-\delta}dX/X$. Derivatives with respect to $s$ introduce fixed
powers of $\log X$ and remain integrable. The interval near zero makes
no contribution. Thus $\mathcal M_W$ is holomorphic in $\Re s>\theta$.

The Mellin transform
\[
 \widehat W(s)=\int_0^\infty W(y)y^s\,\frac{dy}{y}
\]
is entire. For $\Re s>\max(1,\theta)$, absolute convergence of the
Euler product permits interchange of the ideal sum and the integral.
More explicitly, the sum of the integrals of absolute values is
\[
 \left(\int_0^\infty |W(y)|y^{\Re s}\,\frac{dy}{y}\right)
 \sum_{(\mathfrak a,S)=1}
 \frac{|\mu_K(\mathfrak a)\nu(\mathfrak a)|}{(N\mathfrak a)^{\Re s}},
\]
which is finite. The substitution $y=N\mathfrak a/X$ therefore gives
\begin{equation}\label{qr-eq-mellin-identity}
 \mathcal M_W(s)=\widehat W(s)
       \sum_{(\mathfrak a,S)=1}
       \frac{\mu_K(\mathfrak a)\nu(\mathfrak a)}{(N\mathfrak a)^s}
       =\frac{\widehat W(s)}{L_K^S(s,\nu)}.
\end{equation}
The last equality follows by expanding the absolutely convergent
reciprocal Euler product. The identity theorem for meromorphic functions
now gives
\begin{equation}\label{qr-eq-continued-product}
 L_K^S(s,\nu)\mathcal M_W(s)=\widehat W(s)
 \qquad(\Re s>\theta).
\end{equation}
At a pole this is understood as an identity of meromorphic functions.

Suppose that $L_K(\rho,\nu)=0$ and $\Re\rho>\theta$. Since
$\theta\geq0$, the factors deleted in \eqref{qr-eq-deleted-L} are
nonzero at $\rho$, so $L_K^S(\rho,\nu)=0$. Choose a nonzero,
nonnegative function $\phi\in C_c^\infty((1,2);\mathbb R)$ and set
$W(y)=\exp(-\rho\log y)\phi(y)$, where $\log y$ is the real logarithm.
This is an admissible test, and
\[
 \widehat W(\rho)=\int_1^2\phi(y)\,\frac{dy}{y}>0.
\]
But \eqref{qr-eq-continued-product}, evaluated at $\rho$, says that this
positive number is zero. This is a contradiction.
\end{proof}

\begin{remark}\label{qr-rem-mellin-quantifiers}
The constants in \eqref{qr-eq-cancellation-input} may depend on the
fixed character and on the test. The test in the proof is chosen only
after the hypothetical zero is fixed. The theorem excludes an open
half-plane and makes no assertion at its boundary. It does not require
an unsmoothed M\"obius estimate or a bound for $1/L_K$ near a zero.
\end{remark}

For $K=\mathbb Q$, the same proof uses positive integers, a Dirichlet
character $\chi$, and the usual M\"obius function. Thus the bound
\eqref{qr-eq-cancellation-input} with $\theta=1/2$, for every primitive
Dirichlet character and every fixed test, implies GRH by
Theorem~\ref{prim:strip-reduction}. Its functional equation and the
imprimitive Euler factors have already been proved there. The new
arithmetic task would be the square-root estimate itself.

\section{Extracting one sum from a second moment}
\label{qr-sec-moments}

The following elementary lemma makes explicit how the size of a family
affects the exponent obtained for one of its sums.

\begin{lemma}\label{qr-lem-replicated-rows}
Let $(A_u)_{u\in\mathcal R}$ be a finite family of complex numbers,
and let $A$ be another complex number. Suppose there are $J\geq1$
distinct indices $u_1,\ldots,u_J$ such that $|A_{u_j}-A|\leq E$,
where $E\geq0$. Then
\begin{equation}\label{qr-eq-replicated-rows}
 |A|^2\leq\frac2J\sum_{u\in\mathcal R}|A_u|^2+2E^2.
\end{equation}
\end{lemma}

\begin{proof}
For each $j$, the triangle inequality and $(x+y)^2\leq2x^2+2y^2$
give $|A|^2\leq2|A_{u_j}|^2+2E^2$. Average over $j$ and enlarge
the sum of nonnegative terms to all of $\mathcal R$.
\end{proof}

\begin{proposition}\label{qr-prop-power-rows}
Fix an integer $m\geq2$. Suppose, at parameters $D\geq2$ and $H>1$,
that a family satisfies
\[
 \sum_{u\in\mathcal R(H)}|A_u(D)|^2\leq C_\eta D^{1+\eta}H
\]
for an arbitrarily prescribed $\eta>0$. Put $Y=H^{1/m}$ and suppose
that $Y\geq2$. Assume there are at least $cY/\log Y$ distinct indices
in this family at which $A_u(D)$ differs from a distinguished sum
$A_1(D)$ by at most $C D/Y$. Then
\begin{equation}\label{qr-eq-power-extraction}
 |A_1(D)|^2\ll_\eta
 D^{1+\eta}H^{1-1/m}\log H+D^2H^{-2/m}.
\end{equation}
The implied constant may depend on $m,c,C,C_\eta$.
\end{proposition}

\begin{proof}
Apply Lemma~\ref{qr-lem-replicated-rows} with
$J\geq cH^{1/m}/\log(H^{1/m})$ and $E\leq CDH^{-1/m}$.
\end{proof}

\subsection{The sextic example}
\label{qr-subsec-sextic}

We state the arithmetic hypotheses carefully to separate a proved
deduction from its analytic input. Put
$K=\mathbb Q(\sqrt{-3})$, $\omega=e^{2\pi i/3}$, and
$\mathcal O=\mathbb Z[\omega]$. For $z=a+b\omega$ its norm is
$N(z)=|z|^2=a^2-ab+b^2$.
The ring $\mathcal O$ is Euclidean for this norm: writing a complex
number as $x+y\omega$ and rounding $x,y$ to the nearest integers gives
an error $u+v\omega$ with $|u|,|v|\leq1/2$, whose norm is at most
$3/4<1$. Applied to a quotient of two elements, this gives division
with remainder of smaller norm. Hence every ideal is principal.
The six units are $\{\pm1,\pm\omega,\pm\omega^2\}$; their images are
the six distinct units of $\mathcal O/(3)$. Consequently every ideal
coprime to $3$ has a unique generator congruent to $1$ modulo $3$.
We call this its primary generator.

For a prime ideal $\mathfrak p$ not dividing $6$ and an element $u$
not in $\mathfrak p$, the sextic residue symbol
$(u/\mathfrak p)_6$ is the sixth root of unity whose reduction modulo
$\mathfrak p$ equals $u^{(N\mathfrak p-1)/6}$. Set it equal to zero
when $u\in\mathfrak p$, and extend multiplicatively in the ideal
denominator. For the unit ideal the symbol is one for every $u$.
These definitions are possible because the six roots of unity remain
distinct modulo $\mathfrak p$; in particular $6$ divides
$N\mathfrak p-1$.

Fix a finite-order Hecke character $\nu$ and a finite set $S$ containing
the primes above $2$ and $3$ and those dividing its conductor. Define
\begin{equation}\label{qr-eq-sextic-family}
 A_u(D)=\sum_{(\mathfrak n,S)=1}
     \mu_K(\mathfrak n)\nu(\mathfrak n)
     (u/\mathfrak n)_6 W(N\mathfrak n/D).
\end{equation}
The sum $A_1(D)$ is the sum in \eqref{qr-eq-smooth-sum}.

\begin{proposition}[Conditional sextic estimate]\label{qr-prop-sextic}
Assume the following two inputs for this family.
First, for each sufficiently large $Y$, at least $c_S Y/\log Y$
prime ideals outside $S$ satisfy $Y/2<N\mathfrak p\leq Y$, where
$c_S>0$ is fixed. Second, for every fixed $0<\vartheta\leq1/10$,
every $\eta>0$, and every fixed test $W$, one has
\begin{equation}\label{qr-eq-sextic-ms-input}
 \sum_{0<N(u)\leq D^{1+\vartheta}}|A_u(D)|^2
 \ll_{\nu,S,W,\vartheta,\eta}D^{2+\vartheta+\eta}
 \qquad(D\geq2).
\end{equation}
Then for every $\varepsilon>0$ and every fixed test $W$,
\begin{equation}\label{qr-eq-sextic-conclusion}
 A_1(D)\ll_{\nu,S,W,\varepsilon}D^{11/12+\varepsilon}.
\end{equation}
Consequently $L_K(s,\nu)$ has no zero in $\Re s>11/12$.
\end{proposition}

The prime-counting input is the stated consequence of the prime ideal
theorem in this fixed field. Estimate \eqref{qr-eq-sextic-ms-input}
is the principal input asserted in \cite[Section~3]{OpenAIOctober};
that source also states explicit control by a finite number of
derivatives of $W$. Its automorphic proof is not reproduced here.
The following proof establishes the implication from the displayed
inputs, independently of that automorphic proof.
% Source labels: thm:ms, eq:ms, eq:prime-extract.

\begin{proof}
For a prime ideal $\mathfrak p$ outside $S$, let $p$ be its primary
generator. The definition of the residue symbol gives exactly
\[
 (p^6/\mathfrak n)_6=\mathbf1_{\mathfrak p\nmid\mathfrak n}.
\]
Therefore the difference $A_1(D)-A_{p^6}(D)$ contains only ideals
divisible by $\mathfrak p$.

For completeness, the required counting bound is elementary in this
field. The inequality
$a^2-ab+b^2\geq(a^2+b^2)/2$ shows that the number of elements of
$\mathcal O$ of norm at most $T$ is $O(T)$ for $T\geq1$. Since every
ideal is principal, the same upper bound holds for ideals. If $W$ is
supported in $[a,b]$, multiplication by $\mathfrak p$ identifies the
ideals in question with ideals of norm at most $bD/N\mathfrak p$.
If this last quantity is below one there are none; otherwise counting
and $|\mu_K\nu|\leq1$ give, for $Y/2<N\mathfrak p\leq Y$,
\begin{equation}\label{qr-eq-sextic-error}
 |A_1(D)-A_{p^6}(D)|\ll_W D/Y.
\end{equation}

Put $H=D^{1+\vartheta}$ and $Y=H^{1/6}$. The selected elements
$p^6$ have norm at most $H$ and are distinct. Indeed, equality of
sixth powers implies equality of the principal ideals $(p)^6$,
and unique factorization of ideals then implies equality of the
prime ideals. By the assumed prime count there are
$J\gg_S Y/\log Y$ such rows for large $D$.
Lemma~\ref{qr-lem-replicated-rows} gives
\begin{align}
 |A_1(D)|^2
 &\ll_{\nu,S,W,\vartheta,\eta}
 D^{2+\vartheta+\eta}Y^{-1}\log Y+D^2Y^{-2}\notag\\
 &\ll_{\nu,S,W,\vartheta,\eta}
 D^{11/6+5\vartheta/6+\eta}\log D
       +D^{5/3-\vartheta/3}.
 \label{qr-eq-sextic-squared}
\end{align}
For a given $\varepsilon>0$, choose
$\vartheta=\min(1/10,\varepsilon/5)$ and $\eta=\varepsilon/2$.
Use $\log D\ll_\varepsilon D^{\varepsilon/2}$ and take square roots.
The first exponent is at most
$11/12+\varepsilon/12+\varepsilon/2<11/12+\varepsilon$;
the second is at most $5/6$. This proves
\eqref{qr-eq-sextic-conclusion} for large $D$. The bounded remaining
range follows from ideal counting, after increasing the constant.
Theorem~\ref{qr-thm-mellin} gives the zero-free conclusion.
\end{proof}

\subsection{A limitation of this second-moment argument}
\label{qr-subsec-exponent-limit}

\begin{proposition}\label{qr-prop-exponent-limit}
For $m\geq2$ and $\alpha>0$, put
\[
 E_m(\alpha)=\max\left\{
     \frac12+\frac{m-1}{2m}\alpha,\quad
     1-\frac\alpha m\right\}.
\]
Then
\begin{equation}\label{qr-eq-optimal-exponent}
 \min_{\alpha>0}E_m(\alpha)=\frac{m}{m+1},
 \qquad\alpha_{\min}=\frac{m}{m+1}.
\end{equation}
In the restricted range $\alpha\geq1$, the minimum is
$1-1/(2m)$, attained at $\alpha=1$.
\end{proposition}

\begin{proof}
The first affine function is strictly increasing and the second is
strictly decreasing. Equality between them is equivalent to
$\alpha(m+1)=m$. To the left of this intersection the second
function is larger, and to the right the first is larger.
Substitution proves \eqref{qr-eq-optimal-exponent} and uniqueness.
Since $m/(m+1)<1$, the first function is the larger throughout
$\alpha\geq1$; evaluating it at one proves the last assertion.
\end{proof}

Putting $H=D^\alpha$ in \eqref{qr-eq-power-extraction} gives, apart
from an arbitrarily small exponent loss, the bound $D^{E_m(\alpha)}$.
For sextic rows its best exponent over all positive row lengths would
therefore be $6/7$. The actually supplied range
$\alpha=1+\vartheta$, with $\vartheta>0$, has infimum $11/12$.
Thus changing this row length alone cannot obtain exponent $1/2$
from the two terms of \eqref{qr-eq-power-extraction}. This is a
limitation of that particular deduction. It does not disprove a
stronger arithmetic estimate, an alternative extraction, or GRH.
The stronger $7/8$ conclusion stated in \cite{OpenAISeptember}
uses additional estimates; it is not obtained by simply repeating
the extraction above.

\section{Completion and finite iteration}
\label{qr-sec-iteration}

\subsection{Removing added cube factors}

A completed theta expansion can contain additional cube factors.
Their removal is an exact algebraic step before any estimates are made.
The following formulation retains both shared prime factors and
possible zero values of the weights.

\begin{lemma}[M\"obius inversion for cube factors]\label{qr-lem-cube-inversion}
Let $w$ be a function on nonzero integral ideals such that
$w(\mathcal O_K)=1$ and $w(\mathfrak b\mathfrak c)
=w(\mathfrak b)w(\mathfrak c)$ for all ideals. Let $P:(0,\infty)
\to\mathbb C$ satisfy $P(X)=0$ for $0<X<X_0$, for some $X_0>0$.
Define
\begin{equation}\label{qr-eq-cube-completion}
 T(X)=\sum_{\mathfrak b}\frac{w(\mathfrak b)}{N\mathfrak b}
            P\left(\frac{X}{(N\mathfrak b)^3}\right).
\end{equation}
Then both this sum and the following sum are finite, and
\begin{equation}\label{qr-eq-cube-inversion}
 P(X)=\sum_{\mathfrak d}
       \frac{\mu_K(\mathfrak d)w(\mathfrak d)}{N\mathfrak d}
       T\left(\frac{X}{(N\mathfrak d)^3}\right).
\end{equation}
\end{lemma}

\begin{proof}
A nonzero summand in \eqref{qr-eq-cube-completion} must have
$(N\mathfrak b)^3\leq X/X_0$, so only finitely many ideals occur.
Also $T(X)=0$ for $X<X_0$, giving finiteness of the inverse sum.
Substitute \eqref{qr-eq-cube-completion} into its right-hand side
and group by $\mathfrak b=\mathfrak d\mathfrak c$. Complete
multiplicativity and multiplicativity of the norm make the coefficient
of $P(X/(N\mathfrak b)^3)$ equal to
\[
 \frac{w(\mathfrak b)}{N\mathfrak b}
       \sum_{\mathfrak d\mid\mathfrak b}\mu_K(\mathfrak d).
\]
For a nonunit ideal the last sum is the product of one factor
$1-1$ for each distinct prime ideal dividing it, and is zero.
For the unit ideal it is one. Since $w(\mathcal O_K)=1$, the
remaining term is exactly $P(X)$.
\end{proof}

No coprimality between $\mathfrak d$ and $\mathfrak c$ is assumed.
The proof still works when $w$ vanishes at some prime ideals.
In contrast, discarding the added cube terms from a complex sum is
not justified by positivity: they may cancel the other terms.
Splitting \eqref{qr-eq-cube-inversion} at a cutoff leaves a
remainder which must be estimated. This exact inversion is used in
\cite[Section~5]{OpenAIOctober}.
% Source labels: lem:cube-reduction, eq:cube-inverse.

\subsection{A finite descent with an explicit gap}

The next lemma isolates the elementary induction from any analytic
construction of successor sums. The requirement that every successor
remain in the same admissible class is part of its hypotheses.

\begin{lemma}\label{qr-lem-finite-descent}
Fix $B,a,\kappa>0$. For each $D\geq2$, let $\mathcal I(D)$ be a set
of indices. To each index $r$ associate numbers $Q_r\geq0$ and
$0<h_r\leq D^B$, and a set of successors
$\mathcal C(r)\subseteq\mathcal I(D)$. Suppose that
\[
 h_{r'}\leq D^{-a\kappa}h_r\qquad(r'\in\mathcal C(r)).
\]
Assume that for every $\delta>0$ there is $C_\delta\geq1$, independent
of $D$ and of the indices, such that
\begin{align*}
 Q_r&\leq C_\delta D^\delta &&\text{if }h_r\leq1,\\
 Q_r&\leq C_\delta D^\delta
           \left(1+\sup_{r'\in\mathcal C(r)}Q_{r'}\right)
       &&\text{if }h_r>1.
\end{align*}
An empty supremum is zero. Then for every $\varepsilon>0$ there is
$C^*_{\varepsilon,B,a,\kappa}$, independent of $D,r$, such that
$Q_r\leq C^*_{\varepsilon,B,a,\kappa}D^\varepsilon$.
\end{lemma}

\begin{proof}
Put $N=\lceil B/(a\kappa)\rceil$. After $N$ successor steps,
the row parameter is at most $D^{B-Na\kappa}\leq1$.
Fix $\varepsilon>0$, set $\delta=\varepsilon/(N+1)$, and write
$A=C_\delta D^\delta\geq1$. Starting with the terminal bound $A$
and applying the recurrence at most $N$ times gives
\[
 Q_r\leq\sum_{j=1}^{N+1}A^j
      \leq(N+1)C_\delta^{N+1}D^{(N+1)\delta}
      =(N+1)C_\delta^{N+1}D^\varepsilon.
\]
This argument is an induction on the maximum remaining number of
steps. It is valid even with infinitely many successors, because
the inductive bounds are uniform before each supremum is taken.
\end{proof}

In the October argument, the normalized quantity has the form
$Q=\mathcal E(\mathcal H,X,F)/(XF)$ and $\Sigma=XF$. The admissible
range includes $\mathcal H/\Sigma\leq D^{-\kappa}$. Its two analytic
inputs produce successors satisfying
\begin{equation}\label{qr-eq-source-descent}
 \mathcal H'<\mathcal H(\mathcal H/\Sigma)^2
       \leq D^{-2\kappa}\mathcal H,
 \qquad\frac{\mathcal H'}{\Sigma'}\leq
             \frac{\mathcal H}{\Sigma},
 \qquad\Sigma'\leq\Sigma.
\end{equation}
The second and third inequalities preserve these displayed size and gap
conditions, while the first provides the finite descent. Closure of the
other data under each transformation must be established separately. The analytic
proofs of those inputs remain separate obligations. The source
also tracks changing smooth tests: a derivative order $m$ is
replaced by an order at least $4m+12$, and the support interval
enlarges in a controlled way. For a fixed number of steps these
costs can be specified in advance. See
\cite[Section~5]{OpenAIOctober}.
% Source labels: prop:canonical, prop:transfer, sec:canonical-proof.

\begin{proposition}[Fixed parameters cannot be varied without bounds]
\label{qr-prop-nonuniform-kappa}
There exists a nonnegative function $F(D,\kappa)$ such that, for every
fixed $\kappa>0$ and every $\varepsilon>0$,
$F(D,\kappa)=O_{\kappa,\varepsilon}(D^\varepsilon)$ as $D\to\infty$,
but $F(D,1/\log D)=D$ for all $D\geq2$.
\end{proposition}

\begin{proof}
Take $F(D,\kappa)=\exp(1/\kappa)$. At fixed $\kappa$ it is a constant
in $D$, so the first estimate holds with constant $\exp(1/\kappa)$.
Substituting $\kappa=1/\log D$ proves the second assertion.
\end{proof}

Thus a statement proved for every fixed positive gap does not itself
control gaps tending to zero with $D$. In Lemma~\ref{qr-lem-finite-descent}
both the number of steps and the constant may deteriorate as
$\kappa\downarrow0$. To use a growing number of transformations one
must prove the needed quantitative uniformity, including all weight
and derivative costs. Proposition~\ref{qr-prop-nonuniform-kappa}
concerns that logical issue; it is not an assertion about the
actual arithmetic quantities of the source paper.

\section{Uniform continuation for a family of functions}
\label{qr-sec-uniform}

The September paper uses a different analytic arrangement: it
compares two estimates for a sum whose inverse Mellin representation
contains $1/L$. We record a precise abstract version, together with
the uniformity it requires. This is the continuation argument of
\cite[Section~2]{OpenAISeptember}, stated without its arithmetic
construction of the compared sums.
% Source labels: lem:continuation-criterion, eq:uniform-saving.

Let $(L_j)_{j\in\mathcal J}$ be a family of nonzero meromorphic
functions on $\mathbb C$. Suppose none has a zero in $\Re s>1$ and
that $1/L_j$ is bounded on every half-plane $\Re s\geq b>1$,
with the bound allowed to depend on $j,b$. Define
\[
 \beta_*:=\sup\left(\{1/2\}\cup
   \{\Re\rho:\Re\rho\geq1/2,\ L_j(\rho)=0
                        \text{ for some }j\}\right),
\]
where poles are not counted as zeros. Then $1/2\leq\beta_*\leq1$.

\begin{proposition}\label{qr-prop-uniform-continuation}
Fix $\sigma_0\in(1/2,1)$ and suppose $\beta_*>\sigma_0$.
Let $c,b_0\in\mathbb R$ be fixed. Suppose there are numbers
$\omega,\tau$ independent of $j$ with
$0<\omega<\beta_*-\sigma_0$ and $\tau>0$, and, for every $j$,
a holomorphic function $H_j$ on $\Re s>\sigma_0$ such that
$|H_j(s)-1|\leq1/2$ throughout that half-plane. Put
\begin{equation}\label{qr-eq-gaussian-mellin}
 f_j(Z)=\frac1{2\pi i}\int_{\Re s=2}
       Z^{s+c}e^{(s-b_0)^2}\frac{H_j(s)}{L_j(s)}\,ds
       \qquad(Z>0).
\end{equation}
Suppose functions $J_j(Z)$, defined for all sufficiently large $Z$,
satisfy
\begin{align}
 |J_j(Z)|&\ll_j Z^{\sigma_0+c+\omega},
       \label{qr-eq-first-comparison}\\
 |J_j(Z)-f_j(Z)|&\ll_j Z^{\beta_*+c-\tau}.
       \label{qr-eq-second-comparison}
\end{align}
The lower thresholds may depend on $j$. These hypotheses are
inconsistent with $\beta_*>\sigma_0$.
\end{proposition}

\begin{proof}
Let $\delta=\min(\beta_*-\sigma_0-\omega,\tau)>0$.
It is independent of $j$, and $\beta_*-\delta>\sigma_0$.
The triangle inequality gives
\begin{equation}\label{qr-eq-uniform-power-saving}
 |f_j(Z)|\ll_j Z^{\beta_*+c-\delta}
 \quad\text{for all sufficiently large }Z.
\end{equation}

We also need the other endpoint of its Mellin transform. The
integral \eqref{qr-eq-gaussian-mellin} converges absolutely, because
on every fixed vertical strip the Gaussian has modulus a bounded
factor times $e^{-(\Im s)^2}$. The reciprocal and $H_j$ are bounded
on $\Re s\geq2$. For any fixed $B>2$, Cauchy's theorem on rectangles
therefore moves the contour from $2$ to $B$. On the horizontal sides
the Gaussian dominates all remaining factors, so their integrals
tend to zero. Consequently
\begin{equation}\label{qr-eq-small-scale-decay}
 |f_j(Z)|\ll_{j,B}Z^{B+c}\qquad(0<Z\leq1).
\end{equation}
Together \eqref{qr-eq-uniform-power-saving} and
\eqref{qr-eq-small-scale-decay} show that
\[
 F_j(s)=\int_0^\infty f_j(Z)Z^{-s-c}\,\frac{dZ}{Z}
\]
is locally uniformly convergent and holomorphic in
$\Re s>\beta_*-\delta$. Near zero choose $B$ larger than the real
parts on the compact set under consideration; near infinity use
\eqref{qr-eq-uniform-power-saving}. Logarithmic factors introduced
by differentiation are integrable in the same way.

To identify the transform, set
$G_j(s)=e^{(s-b_0)^2}H_j(s)/L_j(s)$ for $\Re s>1$.
Writing $Z=e^u$ in \eqref{qr-eq-gaussian-mellin} gives
\[
 e^{-(2+c)u}f_j(e^u)
   =\frac1{2\pi}\int_{\mathbb R}e^{itu}G_j(2+it)\,dt.
\]
The integrand $G_j(2+it)$ is continuous and integrable. The
left-hand side is integrable in $u$: at $+\infty$ use
$\beta_*\leq1$ and \eqref{qr-eq-uniform-power-saving}, and at
$-\infty$ use \eqref{qr-eq-small-scale-decay} with $B>2$.
Fourier inversion therefore gives
$F_j(2+it)=G_j(2+it)$ for every real $t$.
Both functions are holomorphic on $\Re s>1$, so the identity
theorem gives $F_j(s)=G_j(s)$ throughout that half-plane.

Since $|H_j|\geq1/2$, the function
\[
 R_j(s):=e^{-(s-b_0)^2}\frac{F_j(s)}{H_j(s)}
\]
is a holomorphic continuation of $1/L_j(s)$ to
$\Re s>\beta_*-\delta$. The same $\delta>0$ works for every $j$.
By the definition of the supremum and
$\beta_*-\delta>\sigma_0>1/2$, some $L_j$ has a zero $\rho$ with
$\Re\rho>\beta_*-\delta$. Its reciprocal has a pole there, in
contradiction to this holomorphic continuation. Equivalently, the
meromorphic identity theorem gives $L_j(s)R_j(s)=1$ throughout
the larger half-plane; evaluating at $\rho$ gives $0=1$.
\end{proof}

The constants in the two comparisons need not be uniform in $j$.
The positive saving $\delta$ does need to be uniform. The next
example proves that this distinction cannot be omitted.

\begin{proposition}\label{qr-prop-nonuniform-family}
There is a family of entire functions with zero real parts having
supremum one, although the reciprocal of each member extends
holomorphically from $\Re s>1$ to a strictly larger half-plane.
There is no common positive amount of extension for the whole family.
\end{proposition}

\begin{proof}
For integers $n\geq1$, let
\[
 \rho_n=1-\frac1{n+2},\qquad L_n(s)=s-\rho_n,
 \qquad\delta_n=\frac1{2(n+2)}.
\]
The only zero of $L_n$ is $\rho_n=1-2\delta_n$.
Hence $1/L_n(s)$ is holomorphic on $\Re s>1-\delta_n$,
which extends a positive distance to the left of one for each
fixed $n$. Nevertheless $\rho_n\to1$, so the supremum of their
real parts is one. Given any $\delta>0$, choose $n$ with
$1/(n+2)<\delta$; then $\rho_n>1-\delta$, so the $n$th reciprocal
cannot be holomorphic on $\Re s>1-\delta$.
\end{proof}

These are elementary entire functions, not arithmetic $L$-functions.
Their purpose is to exhibit the quantifier error in replacing a
common positive saving by separate savings which may tend to zero.
In an arithmetic application, the uniform saving and the two
comparisons in Proposition~\ref{qr-prop-uniform-continuation} must
be established by the actual estimates. The continuation argument
does not improve a fixed $11/12$ or $7/8$ half-plane to $1/2$
without those new estimates.

\section{The estimates still needed}
\label{sec:remaining-estimates}

The analytic reductions identify concrete sufficient conditions for GRH.
One is square-root cancellation, with an arbitrarily small power loss,
in the smoothed M\"obius sums of
Theorem~\ref{qr-thm-mellin}, for every fixed primitive character and
every fixed smooth test. The critical-strip reduction then applies.
This is a sufficient condition proved here, not an estimate established
here.

A second route uses the weighted prime sums of the preceding sections.
The announced source constructs upper bounds on long intervals
$[Y,Y^{\lambda}]$ with $\lambda>1$, for suitable finite families of
characters. Three matters must remain explicit: which characters
belong to the family, whether the family is fixed as $Y$ grows, and
whether constants and exponents are uniform in the parameters that
will later vary. A bound only on such intervals cannot be replaced
silently by an eventual bound for every real $x$.

The current source audit locates the input to these bounds in theta
lifts, integral structures, and quantitative bounds for arithmetic
cohomology groups. Their definitions, compatibility conditions, and
parameter ranges have not yet been reduced to a complete checked
ordinary-mathematics proof in this manuscript. No result in the paper
removes this obligation.

The second-moment calculation gives a useful restriction on possible
simplifications. The bound of Proposition~\ref{qr-prop-power-rows},
with its present number of selected indices and its present error,
cannot reach exponent $1/2$ by choosing the length differently.
The best formal exponent for sixth powers is $6/7$ if all positive
length exponents are hypothetically allowed; in the range actually
stated in the October source the limiting exponent is $11/12$.
To reach GRH by that route, an improved arithmetic estimate or a
different use of the family is necessary.

Similarly, the finite descent and the uniform-continuation argument
are useful only with the quantifiers proved in their hypotheses.
A fixed positive gap cannot be replaced by a gap tending to zero
without bounds on the resulting constants and number of steps.
Separate positive improvements for infinitely many characters do not
produce a common improvement. These limitations specify mathematical
work to be done; they are not counterexamples to GRH or to the
announced proof.

\clearpage
\begin{thebibliography}{9}
\small
\raggedright

\bibitem{Alpoge}
L. Alp\"oge, announcement of a proof of the Riemann Hypothesis and
Generalized Riemann Hypothesis, November 2026.
\href{https://alpo.ge/remark.pdf}{Original announcement and links to the released source}.
The announcement credits collaboration with Claude; it does not specify
the model version.

\bibitem{AlpogeSource}
L. Alp\"oge, released formal source accompanying the announcement.
\href{https://github.com/jsmorph/rh/tree/ecd09ff9b4aaf6c08f2b5ef17dc49ca9648a77de}{Public mirror, revision \texttt{ecd09ff9b4aaf6c08f2b5ef17dc49ca9648a77de}}.
The mirror reports an unmodified import of the original archive.
Selected source inspection is distinct from a complete verification.

\bibitem{OpenAISeptember}
OpenAI, \emph{The Quasi-Riemann Hypothesis: A Zero-Free Half-Plane
$\Re(s)>7/8$}, 30 September 2026.
\href{https://github.com/openai/math/blob/adc7f1241b42e322a6451854ab7e4b4c146bf78a/preprints/The-Quasi-Riemann-Hypothesis-September-30-2026/paper.pdf}{Preprint},
\href{https://github.com/openai/math/blob/adc7f1241b42e322a6451854ab7e4b4c146bf78a/preprints/The-Quasi-Riemann-Hypothesis-September-30-2026/build/paper.tex}{TeX source}.
Both links are pinned to revision \texttt{adc7f1241b42e322a6451854ab7e4b4c146bf78a}.

\bibitem{OpenAIOctober}
OpenAI, \emph{The Quasi-Riemann Hypothesis}, 5 October 2026,
written with human assistance.
\href{https://github.com/openai/math/blob/adc7f1241b42e322a6451854ab7e4b4c146bf78a/preprints/The-Quasi-Riemann-Hypothesis-October-5-2026/paper2.pdf}{Preprint},
\href{https://github.com/openai/math/blob/adc7f1241b42e322a6451854ab7e4b4c146bf78a/preprints/The-Quasi-Riemann-Hypothesis-October-5-2026/build/paper2.tex}{TeX source}.
Both links are pinned to revision \texttt{adc7f1241b42e322a6451854ab7e4b4c146bf78a}.

\bibitem{DLMF}
National Institute of Standards and Technology,
\emph{Digital Library of Mathematical Functions},
\href{https://dlmf.nist.gov/1.8}{\S1.8 (Poisson summation)},
\href{https://dlmf.nist.gov/5.2}{\S5.2 (gamma function)}, and
\href{https://dlmf.nist.gov/25.15}{\S25.15 (Dirichlet $L$-functions)}.
Consulted 21 November 2026.

\end{thebibliography}

\end{document}
